%% ========================================================================
%% Preamble: Package dan Konfigurasi
%% ========================================================================

%% ------------------------------------------------------------------------
%% Encoding dan Font
%% ------------------------------------------------------------------------
\usepackage[utf8]{inputenc}      % Input encoding UTF-8
\usepackage[T1]{fontenc}         % Font encoding
\usepackage{lmodern}             % Latin Modern font (lebih baik untuk PDF)
\usepackage{microtype}           % Typographical refinements

%% ------------------------------------------------------------------------
%% Bahasa
%% ------------------------------------------------------------------------
\usepackage[indonesian,english]{babel}  % Dukungan bahasa Indonesia
\selectlanguage{indonesian}

%% Konfigurasi istilah bahasa Indonesia
\addto\captionsindonesian{%
    \renewcommand{\figurename}{Gambar}%
    \renewcommand{\tablename}{Tabel}%
    \renewcommand{\contentsname}{Daftar Isi}%
    \renewcommand{\listfigurename}{Daftar Gambar}%
    \renewcommand{\listtablename}{Daftar Tabel}%
    \renewcommand{\bibname}{Daftar Pustaka}%
    \renewcommand{\indexname}{Indeks}%
    \renewcommand{\appendixname}{Lampiran}%
    \renewcommand{\abstractname}{Abstrak}%
    \renewcommand{\partname}{Bagian}%
    \renewcommand{\chaptername}{Bab}%
}

%% ------------------------------------------------------------------------
%% Geometri Halaman (A4)
%% ------------------------------------------------------------------------
\usepackage{geometry}
\geometry{
    a4paper,
    top=2.5cm,
    bottom=2.5cm,
    left=3cm,
    right=2.5cm,
    bindingoffset=0.5cm,
    headheight=15pt
}

%% ------------------------------------------------------------------------
%% Grafis dan Gambar
%% ------------------------------------------------------------------------
\usepackage{graphicx}            % Untuk memasukkan gambar
\usepackage{float}               % Kontrol posisi float
\usepackage{subfig}              % Subfigure support

%% Konfigurasi caption untuk bahasa Indonesia
\captionsetup{
    figurename=Gambar,
    tablename=Tabel
}

\usepackage{tikz}                % Untuk diagram
\usetikzlibrary{shapes,arrows,positioning,calc}
\usepackage{pgfplots}            % Plot dan grafik
\pgfplotsset{compat=1.18}

%% Path untuk gambar
\graphicspath{{images/}{figures/}}

%% ------------------------------------------------------------------------
%% Tabel
%% ------------------------------------------------------------------------
\usepackage{booktabs}            % Tabel profesional
\usepackage{longtable}           % Tabel panjang (multi-halaman)
\usepackage{tabularx}            % Tabel dengan lebar kolom fleksibel
\usepackage{multirow}            % Sel tabel multi-baris
\usepackage{array}               % Enhanced array and tabular

%% ------------------------------------------------------------------------
%% Warna
%% ------------------------------------------------------------------------
\usepackage{xcolor}              % Dukungan warna
\definecolor{codegreen}{rgb}{0,0.6,0}
\definecolor{codegray}{rgb}{0.5,0.5,0.5}
\definecolor{codepurple}{rgb}{0.58,0,0.82}
\definecolor{backcolour}{rgb}{0.95,0.95,0.92}
\definecolor{linkcolor}{rgb}{0.0,0.2,0.6}
\definecolor{commandcolor}{rgb}{0.1,0.3,0.5}

%% ------------------------------------------------------------------------
%% Listing Kode
%% ------------------------------------------------------------------------
\usepackage{listings}            % Syntax highlighting untuk kode

%% Konfigurasi nama caption untuk listing
\renewcommand{\lstlistingname}{Kode}
\renewcommand{\lstlistlistingname}{Daftar Kode}

\lstdefinestyle{bashstyle}{
    backgroundcolor=\color{backcolour},
    commentstyle=\color{codegreen},
    keywordstyle=\color{magenta}\bfseries,
    numberstyle=\tiny\color{codegray},
    stringstyle=\color{codepurple},
    basicstyle=\ttfamily\small,
    breakatwhitespace=false,
    breaklines=true,
    captionpos=b,
    keepspaces=true,
    numbers=left,
    numbersep=5pt,
    showspaces=false,
    showstringspaces=false,
    showtabs=false,
    tabsize=4,
    frame=single,
    rulecolor=\color{black!30},
    xleftmargin=2em,
    framexleftmargin=1.5em
}

\lstset{style=bashstyle}

%% Definisi bahasa untuk listing
\lstdefinelanguage{bash}{
    keywords={if, then, else, elif, fi, for, while, do, done, case, esac,
              function, return, exit, break, continue, in, select},
    sensitive=true,
    morecomment=[l]{\#},
    morestring=[b]",
    morestring=[b]'
}

%% ------------------------------------------------------------------------
%% Hyperlinks dan Referensi
%% ------------------------------------------------------------------------
\usepackage[
    colorlinks=true,
    linkcolor=linkcolor,
    urlcolor=blue,
    citecolor=linkcolor,
    bookmarks=true,
    bookmarksnumbered=true,
    pdfborder={0 0 0}
]{hyperref}

\usepackage[capitalise,noabbrev]{cleveref}  % Referensi pintar

%% Konfigurasi cleveref untuk bahasa Indonesia
\crefname{figure}{gambar}{gambar}
\Crefname{figure}{Gambar}{Gambar}
\crefname{table}{tabel}{tabel}
\Crefname{table}{Tabel}{Tabel}
\crefname{chapter}{bab}{bab}
\Crefname{chapter}{Bab}{Bab}
\crefname{section}{bagian}{bagian}
\Crefname{section}{Bagian}{Bagian}
\crefname{subsection}{subbagian}{subbagian}
\Crefname{subsection}{Subbagian}{Subbagian}
\crefname{equation}{persamaan}{persamaan}
\Crefname{equation}{Persamaan}{Persamaan}
\crefname{listing}{kode}{kode}
\Crefname{listing}{Kode}{Kode}
\crefname{appendix}{lampiran}{lampiran}
\Crefname{appendix}{Lampiran}{Lampiran}

%% ------------------------------------------------------------------------
%% Header dan Footer
%% ------------------------------------------------------------------------
\usepackage{scrlayer-scrpage}    % KOMA-Script header/footer
\clearpairofpagestyles
\ihead{\leftmark}
\ohead{\rightmark}
\cfoot[\pagemark]{\pagemark}
\pagestyle{scrheadings}

%% ------------------------------------------------------------------------
%% Bibliography (Daftar Pustaka)
%% ------------------------------------------------------------------------
\usepackage[
    backend=biber,
    style=ieee,
    sorting=none
]{biblatex}
\addbibresource{references.bib}

%% ------------------------------------------------------------------------
%% Matematika
%% ------------------------------------------------------------------------
\usepackage{amsmath}             % Matematika lanjutan
\usepackage{amssymb}             % Simbol matematika
\usepackage{amsthm}              % Theorem environments

%% ------------------------------------------------------------------------
%% Box dan Environment Khusus
%% ------------------------------------------------------------------------
\usepackage[most]{tcolorbox}     % Colored boxes

% Box untuk catatan
\newtcolorbox{notebox}{
    colback=blue!5!white,
    colframe=blue!75!black,
    title=Catatan,
    fonttitle=\bfseries
}

% Box untuk peringatan
\newtcolorbox{warningbox}{
    colback=red!5!white,
    colframe=red!75!black,
    title=Peringatan,
    fonttitle=\bfseries
}

% Box untuk tip
\newtcolorbox{tipbox}{
    colback=green!5!white,
    colframe=green!75!black,
    title=Tips,
    fonttitle=\bfseries
}

% Box untuk contoh
\newtcolorbox{examplebox}[1][]{
    colback=yellow!10!white,
    colframe=orange!75!black,
    title=Contoh: #1,
    fonttitle=\bfseries
}

% Box untuk latihan
\newtcolorbox{exercisebox}[1][]{
    colback=purple!5!white,
    colframe=purple!75!black,
    title=Latihan #1,
    fonttitle=\bfseries
}

%% ------------------------------------------------------------------------
%% Command Khusus Linux
%% ------------------------------------------------------------------------
\newcommand{\cmd}[1]{\texttt{\textcolor{commandcolor}{#1}}}
\newcommand{\file}[1]{\texttt{#1}}
\newcommand{\dir}[1]{\texttt{#1/}}
\newcommand{\var}[1]{\texttt{\$#1}}
\newcommand{\keystroke}[1]{\framebox{\texttt{#1}}}

%% ------------------------------------------------------------------------
%% Indeks
%% ------------------------------------------------------------------------
\usepackage{makeidx}
\makeindex

%% ------------------------------------------------------------------------
%% Teorema dan Definisi
%% ------------------------------------------------------------------------
\theoremstyle{definition}
\newtheorem{definition}{Definisi}[chapter]
\newtheorem{example}{Contoh}[chapter]

\theoremstyle{plain}
\newtheorem{theorem}{Teorema}[chapter]

%% Konfigurasi nama untuk cleveref (teorema)
\crefname{definition}{definisi}{definisi}
\Crefname{definition}{Definisi}{Definisi}
\crefname{example}{contoh}{contoh}
\Crefname{example}{Contoh}{Contoh}
\crefname{theorem}{teorema}{teorema}
\Crefname{theorem}{Teorema}{Teorema}

%% ------------------------------------------------------------------------
%% Spacing dan Format
%% ------------------------------------------------------------------------
\usepackage{setspace}
\onehalfspacing                  % 1.5 line spacing

\setlength{\parindent}{1.5em}    % Indentasi paragraf
\setlength{\parskip}{0.5em}      % Jarak antar paragraf

%% ------------------------------------------------------------------------
%% Chapter dan Section Formatting
%% ------------------------------------------------------------------------
\addtokomafont{chapter}{\Huge}
\addtokomafont{section}{\LARGE}
\addtokomafont{subsection}{\Large}

%% ------------------------------------------------------------------------
%% Miscellaneous
%% ------------------------------------------------------------------------
\usepackage{lipsum}              % Dummy text (hapus untuk versi final)
\usepackage{enumitem}            % Enhanced list environments
\usepackage{csquotes}            % Context sensitive quotation

%% Suppress overfull hbox warnings di batas tertentu
\hfuzz=5pt
